\chapter{Introduction}

Autonomous navigation systems have traditionally been designed around the perspective of the moving robot itself. Cameras and sensors are mounted directly on the platform to capture a ground-level, ego-centric view of the environment, mirroring how a human driver perceives the road ahead. This design works well in environments where the navigation path is physically distinguishable: lane markings, boundary edges, or clear contrast between the drivable surface and its surroundings give the system enough visual information to follow a route. However, this fundamental assumption breaks down in a wide class of practical applications where no such ground path exists.

Consider autonomous lawn mowers operating across uniform grass, or robotic vacuum cleaners navigating featureless carpet and bare flooring. In these environments, the ground surface provides no inherent visual guidance given that there are no lane markings to follow, no edges to track, and no color contrast to exploit. Conventional ego-centric navigation approaches, including modern end-to-end deep learning systems trained to recognize road features, are architecturally unsuited to these conditions without significant redesign. A solution must be explored that does not depend on the robot's ability to see and recognize a physical ground path from its own perspective.

Beyond the absence of a visible path, existing autonomous navigation frameworks impose further constraints that limit their deployment in low-cost and resource-constrained settings. Existing solutions like Simultaneous Localization and Mapping (SLAM) rely on expensive sensor arrays like LiDAR, depth cameras, and inertial measurement units which is then paired with computationally intensive software pipelines that demand high-performance onboard processors. For compact consumer-grade platforms such as robotic vacuum cleaners and autonomous lawn mowers, these requirements are fundamentally mismatched with the operational realities of tight unit cost, limited power budgets, and minimal onboard processing capacity.

This study proposes an alternative navigation framework that addresses both of these constraints simultaneously. By separating path perception from the robot's onboard egocentric view, the system offloads visual processing to a fixed overhead camera and employs planar ArUco fiducial markers as geometric reference anchors. A homography transformation converts the overhead camera feed into a rectified top-down representation of the environment, and a lightweight Convolutional Neural Network trained through behavioral cloning maps this geometrically simplified visual input directly to continuous steering commands without any classical controller feedback, without expensive sensors, and without requiring a physically marked ground path. This establishes a foundational proof-of-concept for autonomous navigation in environments where the ground surface provides no inherent visual guidance, deployable on resource-constrained hardware.
\section{Background of the Study}

Autonomous navigation is the ability of robots and autonomous vehicles to generate a path and move along it without human intervention, utilizing data captured from onboard sensors and sometimes remote navigation aids \cite{Cebollada2020MobileRoboticsAI}. These systems consistently rely on simultaneous localization and mapping (SLAM) for autonomous navigation. Additionally, these systems also utilize LiDAR, depth cameras, or multi-sensor fusion \cite{Liu2024IndoorSensingReview, Wei2024SLAMResearchReview, Zhu2023CameraLidarIMUFusionSLAMSurvey}. While generally effective, these robotic systems demand high computational costs, expensive robotics components, and complex software integration requirements, which makes them less accessible for developing low-cost robotics platforms \cite{Chea2023CollaborativeConstructionRobot}. To address these limitations, previous studies have already explored computer-vision navigation using fiducial markers, which offer a simple, lightweight, and reliable framework for localization, particularly in structured environments \cite{Alghamdi2025MobileRobotMarkersMapping}.

According to \cite{kuzmenko2025fiducial}, a fiducial marker refers to a specialized code-type marker used in computer vision, which is set in an environment or scene to help imaging systems find reference points. The robustness of fiducial markers such as ArUco and AprilTag in providing low-cost localization has already been explored in a 2025 systematic review of the mapping of mobile robot navigation \cite{Alghamdi2025MobileRobotMarkersMapping}. The study of \cite{Mantha2022FiducialNetworkNavigation} concluded that marker size, spacing, and the placement of a fiducial marker may affect robot navigation, well-designed markers show potential in achieving high success rates with minimal computational overhead. 

Several recent studies probed the ability of ArUco markers to provide indoor localization for autonomous robot navigation. Research on planar indoor localization using ArUco \cite{Hinderer2025ArucoPlacementIndoor} and multi-marker setups illustrates that planar marker fields can deliver accurate 2-D pose estimates which are appropriate for autonomous robot navigation \cite{Lin2023MarkerBasedLocalization}, especially when combined with simple filters or odometry. To further maximize localization, multi-camera setups are also explored, as seen in the study of \cite{Sevgi2025MulticameraIndoorAruco}. Studies on fiducial‑marker networks for indoor navigation \cite{Braun2022RobotAtFactoryLocalization} explicitly frame markers as an alternative to the expensive SLAM for building service and inspection robots, highlighting low instrumentation and computation as key advantages.

For robot navigation, recent research has shown interest in using end-to-end frameworks that map visual input directly to steering \cite{Han2024NeuPAN} or obstacle avoidance actions, usually with deep reinforcement learning or attention-based architecture \cite{Zhou2025DRLSpatiotemporalTransformer, Yu2025RLAUVMotionSystems}. Although such methods can achieve strong obstacle avoidance performance in simulation or structured test environments, multiple studies emphasize issues with stability, overfitting, and susceptibility to noise, leading to poor generalization outside training conditions. 

Despite the demonstrated effectiveness of ArUco markers for localization and the growing body of work on end-to-end visual navigation, relatively few studies have explored combining these two domains. Most marker-based systems stop at pose estimation and hand off to classical controllers, while most end-to-end systems operate on raw ego-centric feeds without geometric preprocessing. This study addresses that intersection directly.

\section{Statement of the Problem} 
Current ego-centric, vision-based autonomous navigation models rely heavily on the presence of physical path features, such as lane markings, boundary edges, or distinct color contrasts between the drivable surface and its surroundings. When deployed in featureless environments such as uniform grass, carpet, or bare flooring, these foundational assumptions break down. Without a visually distinguishable ground path, conventional end-to-end deep learning approaches fail to generate reliable steering commands, rendering them fundamentally inapplicable to these environments without significant architectural redesign or the reliance on expensive, computationally heavy SLAM sensor arrays.

Conversely, while alternative localization methods using fiducial markers such as ArUco have proven effective for low-cost indoor tracking, existing marker-based systems predominantly hand off pose estimation or localization to classical control algorithms rather than utilizing neural network integration. There is a distinct gap in the literature regarding the intersection of fiducial marker localization and end-to-end imitation learning. Specifically, it remains unproven whether offloading visual perception to a fixed overhead camera and applying planar homography can simplify visual inputs sufficiently to allow a lightweight Convolutional Neural Network (CNN) to map those inputs directly to continuous steering commands, bypassing the need for classical controllers and physical ground paths entirely.

\section{Research Objective} 
The following are the general objective and specific objectives of the study.

\subsection{General Objective}
To develop an autonomous vehicle system with external top-down camera and utilizing an end-to-end Convolutional Neural Network (CNN) architecture that maps raw visual input, digitally augmented with a virtual ground path, directly to robot’s steering.
path.

\subsection{Specific Objectives}
\begin{enumerate}
    \item To implement a robust Inverse Perspective Mapping (IPM) pipeline using ArUco markers for dynamic homography calibration, capable of transforming angled, ego-centric RGB camera frames into rectified, top-down, bird’s-eye view (BEV) representations of the local path.
    \item To construct a specialized behavioral cloning dataset by manually driving the robot, consisting of synchronized pairs of preprocessed BEV path images and human-demonstrated steering angles, collected across various track configurations.
    \item To design and train a regression-based Convolutional Neural Network (CNN) architecture that maps high-dimensional BEV visual inputs directly to continuous steering commands (and optionally linear velocity), utilizing Mean Squared Error (MSE) loss for optimization.
    \item To deploy the trained model onto the robot's driving pipeline such that inference output is sent directly to the motor drivers, enabling fully autonomous navigation.
    \item To evaluate the autonomous robot by measuring lane‑keeping accuracy (CTE/RMSE), heading error, trial‑to‑trial consistency, and steering actuation behaviour, and summarizing results using confusion matrices and per‑scenario statistics.
\end{enumerate}

%\section{Proposed Solution} 

\section{Scope and Limitations of the Research}
This study focuses on the development of a low-cost, end-to-end autonomous navigation system using a single overhead RGB camera. The system utilizes Inverse Perspective Mapping (IPM) facilitated by ArUco markers to process visual inputs.

The hardware implementation is restricted to computationally constrained platforms (e.g., Raspberry Pi or NVIDIA Jetson Nano) to validate the lightweight nature of the proposed architecture. The robot operates exclusively in structured environments where ArUco markers can be reliably detected, and no 3-D mapping, LiDAR, or depth-based perception is included. Unlike traditional SLAM systems that integrate probabilistic mapping and multi-sensor fusion \cite{Cadena2016SLAMRobustPerception, Cieslewski2017DecentralizedVisualSLAM}, this study does not attempt full environment reconstruction or dynamic obstacle mapping. Instead, localization is constrained to planar marker fields consistent with research demonstrating accurate 2-D pose estimation using fiducials \cite{GarridoJurado2014FiducialMarkers, RomeroRamirez2018SpeededFiducialDetection}.

The study is also limited to a differential drive (2WD) chassis. This is intentional as the Convolutional Neural Network is not trained to output complex holonomic movements or deal with Ackermann steering physics.

As for the environmental scope of this study, the navigation environment is limited to structured, planar indoor surfaces with consistent artificial lighting. The robot operates within an area of interest, which is a predefined 2-D plane where ArUco markers are visible and explicitly placed to define track boundaries or reference points. The system is not designed for outdoor environments, uneven terrain, or varying illumination conditions.

Neural-network steering prediction is trained only on data collected within the experimental environment. Thus, the learned model is not expected to generalize to unstructured or markerless environments, consistent with known limitations of end-to-end visual navigation models regarding generalization and noise sensitivity \cite{Bansal2020ChauffeurNet, Codevilla2019BehaviorCloningLimits}. The evaluation metrics, cross-track accuracy, heading error, and success rate, are limited to controlled indoor trials and do not include long-term autonomy, multi-robot coordination, or obstacle-rich navigation.

\section{Significance of the Research } 
This study contributes to the growing demand for accessible and low cost autonomous driving robot navigation frameworks by presenting a system that removes the need for expensive sensors, advanced SLAM algorithms, and high-end computing hardware. While previous works have highlighted that fiducial markers such as Aruco provide efficient localization \cite{GarridoJurado2014FiducialMarkers, RomeroRamirez2018SpeededFiducialDetection}, relatively few studies have used these markers into a complete low cost 2-D navigation system suited for educational or research purposes.

This research contributes to the field of robotic control by exploring the potential of End-to-End Imitation Learning in constrained environments. While previous works have validated ArUco markers for localization \cite{RomeroRamirez2018SpeededFiducialDetection} and CNNs for visual navigation separately, this study bridges these domains. It investigates whether a homography-based pre-processing step can simplify visual inputs enough to allow a lightweight neural network to provide navigation controls accurately without the computational overhead of deeper, complex architectures.

The study is locally and globally significant because it provides:

\begin{enumerate}
    \item A framework for implementing advanced Deep Learning concepts using only a standard webcam and a microcontroller, making the technology accessible for educational and low-resource industrial applications.
    \item A simple experimental platform for teaching concepts related to localization, control systems, and machine learning.
    \item A replicable and scalable framework for researchers exploring low-cost autonomous navigation in structured indoor spaces.
\end{enumerate}

Overall, this study aims to address the gaps in the literature by producing a fully integrated, low-cost navigation framework that combines fiducial-marker localization and End-to-End robot control, offering a practical solution for accessible autonomous robot navigation systems.

%\blindtext
